% 02 · The data — surveys, abundance space, kinematic truth, the cluster set
\section{THE DATA}
\label{sec:data}

\deckslides{slides 2 (the data) and 21 (the scale problem).}

Three surveys supply everything this workbook measures. APOGEE provides the
chemistry, Gaia the astrometry and radial velocities, and GALAH the southern
main sequence that appears in \S\ref{sec:physics}. The samples are built from
public data releases, and the whole pipeline --- download, quality cuts,
membership, matrices --- is reproducible with three commands
(\S\ref{sec:howtorun}).

\subsection{APOGEE: abundances for giants}

The Apache Point Observatory Galactic Evolution Experiment
\citep{Majewski:17} is a near-infrared
($H$-band, $1.51$--$1.70$ $\mu$m) spectroscopic survey of $\sim\!10^5$ stars
per data release, obtained with the 2.5-m Sloan Foundation Telescope in the
north and the du Pont Telescope in the south. The infrared window matters for
this project: it sees through the dust that hides the inner disc, and it
targets cool giants, which are bright enough to be observed across the whole
Galaxy.

Until recently the abundances came from the ASPCAP pipeline. The SDSS-V
release used here --- DR19 \texttt{astraAllStarASPCAP} \citep{Almeida:23},
$1.17$ GB, $1\,095\,480$ targets --- is produced by its successor, the
\emph{Astra} framework. The transfer is not cosmetic; two schema changes are
worth knowing about because they bite silently rather than loudly:

\begin{itemize}
\item abundances are published as \emph{$[$X/H$]$}, not $[$X/Fe$]$ under a
  different name, so element ratios must be formed explicitly as
  $[$X/Fe$] = [$X/H$] - [$Fe/H$]$;
\item the quality flags were renamed --- \texttt{ASPCAPFLAG} became
  \texttt{flag\_bad}, \texttt{STARFLAG} became \texttt{spectrum\_flags}.
\end{itemize}

DR19 also reanalyses rather than replaces DR17: 717\,689 of its rows carry
\texttt{release='dr17'}, all reprocessed by a single version of the pipeline.
That detail becomes important in \textbf{Section~\ref{sec:spectral}}, where an
earlier version of our benchmark was nearly derailed by mixing data
\emph{products} from the two.

\begin{marginnote}[]
\entry{ASPCAP / Astra}{The abundance pipelines: ASPCAP up to SDSS-IV DR17,
Astra from SDSS-V DR19 onwards. Same philosophy --- fit synthetic spectra
to the observation --- different catalogues and column names.}
\entry{Abundance precision}{The per-element statistical error on $[$X/H$]$,
typically $0.02$--$0.1$ dex for APOGEE giants, depending on the element and
the signal-to-noise ratio. It sets the ceiling on what chemistry can tag.}
\end{marginnote}

\subsection{Gaia and GALAH: the referee and the main sequence}

Gaia DR3 \citep{Gaia:23} supplies parallaxes, proper motions and the
cross-matched photometry for every star in the sample, all of it carried
inside the APOGEE catalogue for convenience. These are the quantities that
define membership (\S\ref{sec:kinematics}); we never feed them to a
tagging method, and \textbf{Section~\ref{sec:validation}} explains why that
asymmetry is the whole point.

GALAH DR4 \citep{DeSilva:15, Buder:21} covers declination $\lesssim +25^\circ$
with a high-resolution optical spectrograph, which means it sees the
\emph{main sequence} where APOGEE sees giants. It enters only in
\S\ref{sec:physics}, where the two surveys are combined to recover a cluster
age and distance at the same time.

\subsection{The abundance space}

Everything is standardised into one matrix: 16 numbers per star, one column
per element, built from the survey catalogue by a fixed sequence of steps.
Each step is a modelling choice, not housekeeping, and each has a failure mode
we have hit.

\begin{enumerate}
\item \textbf{Quality cuts.} Signal-to-noise ratio $\geq 100$, ASPCAP flag
  zero, stellar flag zero. Cuts this aggressive remove the noisy tail at the
  cost of sample size; loosening them changes the answer
  (\S\ref{sec:benchmark}, the levers).
\item \textbf{The element set.} C, N, O, Na, Mg, Al, Si, S, K, Ca, Ti, V,
  Cr, Mn, Ni --- the $\alpha$ and iron-peak elements --- plus $[$Fe/H$]$ as
  the sixteenth coordinate. This is the set used by the reference papers, so
  results are comparable with them.
\item \textbf{Missing values.} Weak lines go undetected in metal-poor stars
  and their abundances come back as NaN. Dropping every star with a NaN
  deletes the metal-poor globulars \emph{entirely} --- M~15 and M~92 vanish
  from the sample, and with them the cleanest tagging cases in the whole
  project. We therefore keep stars with at least 8 of 16 finite elements and
  fill the remaining gaps with the column median. This single choice is the
  difference between a sample that contains the interesting objects and one
  that does not.
\item \textbf{Standardisation.} Each element is shifted to zero median and
  scaled to unit standard deviation, so that no element dominates the distance
  metric by virtue of its units. (An early bug in this project applied
  inverse-variance weights \emph{before} this step, which silently cancelled
  them; see \S\ref{sec:benchmark}.)
\item \textbf{Row normalisation.} Each star's 16-vector is scaled to unit
  $L_2$ norm. The reason is metric compatibility: after this step, Euclidean
  distance is a monotone function of cosine distance, so t-SNE and UMAP (which
  use Euclidean) and EVoC (whose native metric is cosine) all see the same
  geometry, and the comparison between them is fair. It is also the single
  most effective precision lever in the pipeline: without it, density-based
  clustering merges the whole field into one blob.
\end{enumerate}

\begin{figure}[htbp]\centering
  \fig[0.62]{cspace.png}
  \caption{The shape of the problem, in two of the sixteen dimensions (from
  the deck). Field stars in grey; the members of one cluster in red. The
  cluster is a tight, slightly elongated clump embedded in a cloud several
  orders of magnitude larger --- a small overdensity in a big, structured
  background. Everything in
  \S\ref{sec:fundamentals}--\S\ref{sec:evoc} is a strategy for finding
  objects like this one.}
  \label{fig:cspace}
\end{figure}

\subsection{Kinematic membership is the referee, not a feature}
\label{sec:kinematics}

To score anything we need labels, and this is where chemical tagging is
unusually lucky: a cluster is also a kinematic object. Its members currently
share a position on the sky, a parallax, a proper motion and a radial
velocity, because they were born at the same place with almost the same
velocity. We use those four observables, $\sigma$-clipped around the cluster
centre, to label each star \emph{member} or \emph{field}.

The circularity risk is obvious: if kinematics are used to build the labels,
then a method that receives kinematics as an input is being scored on its own
input. We therefore treat kinematics as a \emph{ceiling}, not a competing
signal --- the number that says how well \emph{any} chemistry-based method
could possibly do on this sample. Published cluster catalogues provide an
independent cross-check: scoring against Simbad memberships gives numbers
within a few percent of the kinematic ones, and both are reported in the
benchmark \citep{Dias:02, Harris:96}.

\subsection{The cluster set}
\label{sec:clusters}

The target list is the 23 clusters of \citet{GarciaDias:19} plus the Pleiades
of \citet{Kos:17} --- and, for the two-survey work, the southern open clusters
NGC~2243 and Collinder~261. Twenty-five objects in total: seven globular and
eighteen open. They span two orders of magnitude in age, from the Pleiades at
$\sim\!100$ Myr to NGC~188 at $7.6$ Gyr, and two orders of magnitude in mass.
The counts below come from the full DR19 run (357\,056 field stars). They are
\emph{row counts}, not star counts: the sample contains repeat entries for the
same star, a defect quantified in \S\ref{sec:validation} and one we flag here
rather than at the end because it affects how every macro average in this
workbook should be read.

\begin{table}[htbp]
\tabcolsep5pt
\caption{The 25 target clusters and the number of kinematic members each
contributes in the DR19 sample. Globulars are marked with a $\dagger$. The
asymmetry in the counts is real and consequential: M~67 alone supplies
$23\%$ of all members, while three clusters supply fewer than ten.}
\label{tab:clusters}
\begin{center}
\begin{tabular}{@{}lrlr@{}}
\hline
Cluster & $n_{\rm mem}$ & Cluster & $n_{\rm mem}$\\
\hline
Berkeley 17 & 7 & M 3$^{\dagger}$ & 154\\
Berkeley 66 & 12 & M 5$^{\dagger}$ & 67\\
Berkeley 71 & 14 & M 13$^{\dagger}$ & 34\\
Collinder 261 & 22 & M 15$^{\dagger}$ & 31\\
IC 166 & 35 & M 71$^{\dagger}$ & 39\\
King 5 & 12 & M 92$^{\dagger}$ & 25\\
King 7 & 8 & NGC 1245 & 21\\
M 67 & 230 & NGC 1798 & 13\\
Pleiades & 23 & NGC 188 & 25\\
& & NGC 2158 & 6\\
& & NGC 2243 & 36\\
& & NGC 2420 & 19\\
& & NGC 6791 & 45\\
& & NGC 6819 & 62\\
& & NGC 7789 & 47\\
\hline
\multicolumn{4}{@{}l}{M 107$^{\dagger}$ contributes 15 members.\quad
Total: 1\,002 members in 25 clusters.}\\
\hline
\end{tabular}
\end{center}
\begin{tabnote}
Counts from the full-sky DR19 run (\texttt{cluster run}), field star cap
disabled, on the member frame the DR19 baseline scores --- every member row
with at least nine finite abundances out of sixteen. Reproduce with
\texttt{python scripts/population\_counts.py --baseline-min-finite 9};
the counts are identical for any floor between nine and fifteen. Member counts
differ slightly between runs because the clustering must first find a group
before its overlap with the truth can be measured; these are the truth-side
counts.
\end{tabnote}
\end{table}

\begin{figure}[htbp]\centering
  \begin{tabular}{@{}c@{\hskip8pt}c@{}}
    \fig[0.42]{sky_m67.png} &
    \fig[0.42]{sky_m3.png}
  \end{tabular}
  \caption{Two of the targets, in the optical (DSS2 colour composites centred
  on the cluster): M~67 (left), an old solar-metallicity open cluster, and
  M~3 (right), a metal-poor globular. The two objects illustrate the range of
  the sample and, in the results, the range of difficulty: the globular is
  chemically distinct from the field around it, the open cluster is not.}
  \label{fig:skysamples}
\end{figure}

\subsection{How to reproduce the sample}
\label{sec:howtorun}

The pipeline is one package with a single configuration file, so every number
in this workbook can be regenerated from a fresh checkout. It is written on
\texttt{numpy} \citep{Harris:20}, \texttt{pandas} \citep{McKinney:10},
\texttt{scikit-learn} \citep{Pedregosa:11}, \texttt{astropy}
\citep{Astropy:22} and \texttt{matplotlib} \citep{Hunter:07}, with the
clusterers taken from their reference implementations. The commands below
are all you need to rebuild the sample and the benchmark.

\begin{summary}[HOW WE RAN IT]
\begin{quote}\small
\begin{verbatim}
git clone github.com/garciadias/iaa-advanced-neural-networks-2026-draft
cd iaa-advanced-neural-networks-2026-draft
uv sync                       # Python 3.13 environment
uv run cluster download       # 1.17 GB, SDSS-V DR19 Astra ASPCAP
uv run cluster run --fast     # ~1 min on a laptop, capped field
CLUSTER_FAST=0 uv run cluster run   # full all-sky run, 10-20 min
\end{verbatim}
\end{quote}
Every flag lives in \texttt{src/cluster/config.py} and every flag also has an
environment override with the \texttt{CLUSTER\_} prefix, which is how the
variants in \S\ref{sec:benchmark} were produced. The Docker path
(\texttt{docs/docker.md}) exists for machines without \texttt{uv}.
\end{summary}

\begin{table}[htbp]
\centering
\small
\tabcolsep4pt
\begin{tabular}{@{}ll@{}}
\toprule
setting & value used throughout \\
\midrule
\texttt{SNR\_MIN} & $100$ \\
dwarf cut (\texttt{DWARF\_ONLY}) & off \\
elements & 16: \texttt{[C/Fe]} \dots\ \texttt{[Ni/Fe]} (15 ratios) $+$ \texttt{[Fe/H]} \\
imputation & median, requiring $\geq\!8$ finite elements \\
standardisation & on (zero median, unit s.d.) \\
row normalisation & on \\
element weights & off \\
\texttt{HDBSCAN.min\_cluster\_size} & $5$ \\
\texttt{TSNE.perplexity} / \texttt{UMAP.n\_neighbors} & $30$ / $15$ \\
\texttt{random\_state} & $42$ \\
\bottomrule
\end{tabular}
\caption{The pipeline's settings for every number in this workbook, unless a
variant is named. Overrides are environment variables with the
\texttt{CLUSTER\_} prefix --- \texttt{CLUSTER\_NORMALIZE\_ROWS=0} is the one
that most changes the answer.}
\label{tab:config}
\end{table}

\begin{issues}[EXERCISES]
\begin{enumerate}
\item Take the 16-element vector for one star in the sample. Standardise it
  and then row-normalise it. Show algebraically that after row normalisation
  the Euclidean distance between two stars is a strictly monotone function of
  the cosine distance between them. Why does EVoC care?
\item \texttt{MIN\_FINITE\_ELEMENTS} defaults to 8 of 16. Set it to 16 (strict
  complete-case) and count how many of the 25 clusters survive with at least
  five members. Which clusters are lost, and what does that do to any
  conclusion about the surviving ones?
\item M~67 contributes 230 of 1\,002 members. Recompute a macro-average
  metric over clusters with and without M~67 in the pool. How much of the
  published summary depends on one cluster?
\end{enumerate}
\end{issues}
